
\section{Conclusion}
\label{sec:conclusion}

In this paper, we proposed ApproxPilot-SE, an automated approximate accelerator optimization framework. The framework employs configuration-aware GNNs to predict hardware PPA and application quality, incorporating path-aware aggregation and measured-reference local calibration for Delay estimation. By integrating search-guided measurements into iterative surrogate updates, ApproxPilot-SE enables the surrogate models and DSE process to evolve together, improving the utilization of available synthesis data.

Experiments on five domain-specific accelerators, including DCT and Laplace with substantially larger combinatorial design spaces, demonstrate the effectiveness of ApproxPilot-SE. On DCT, the proposed framework achieves 3.79 times PPA HV Gain of the strongest baseline under the same training budget at 90\% SSIM retention. Notably, even with 40\% fewer training labels, it achieves 3.42 times the gain of the strongest baseline using the full budget. ApproxPilot-SE also maintains superior optimization performance under stringent accuracy constraints, achieving the highest HV Gain across all five benchmarks at 99\% SSIM retention.


A representative DCT design preserves 99.62\% of the original SSIM while reducing area, power, and delay by 19.36\%, 20.68\%, and 9.70\%, respectively. End-to-end ablation studies further validate the effectiveness of iterative feedback, guided candidate selection, and latency modeling in improving the final optimization results. Together, these findings demonstrate the capability of ApproxPilot-SE to efficiently explore large approximate accelerator design spaces and achieve superior PPA--accuracy tradeoffs across different applications.

